% ===============================================================
% Abstract  --  Integration and Reporting Lead
% ===============================================================
This report presents a Shakespeare-aware language system that helps a
user with no prior exposure to Shakespeare ask questions about
\emph{Hamlet}, \emph{Macbeth}, and \emph{Romeo and Juliet} and receive
accurate, evidence-grounded explanations. We treat the task as one of
domain adaptation: rather than pretraining or full fine-tuning, we adapt
a compact pretrained model, Gemma~3~4B running locally via Ollama, using
retrieval-augmented generation (RAG) and prompt design. The provided
structured corpus is prepared at two granularities, scene-level chunks
and speaker-level utterances, with each retrieval unit enriched by a
modern-English scene summary and keywords and indexed by a MiniLM
sentence embedding. Three systems are compared: a no-retrieval baseline,
a standard RAG pipeline, and an enhanced pipeline that adds
cross-encoder reranking and scene-level diversity. The system supports
four interaction modes, concept explanation, contextual question
answering, evidence retrieval, and clearly labelled stylised generation,
and always displays retrieved evidence alongside its answers. We
evaluate on fifteen questions (five instructor-provided, ten
group-designed) using a five-criterion rubric. Retrieval relevance is
consistently strong, while the analysis isolates the generation stage as
the primary bottleneck and the main target for future improvement.
